\cvsection{Education}

\begin{cventries}

  \cventry
    {Instituto Superior Técnico - Universidade de Lisboa}
    {PhD. Electrical \& Computer Engineering}
    {Lisbon, Portugal}
    {2026 -- Present}
    {
      \begin{cvitems}
        \item Researching low-overhead event-based performance monitoring architectures for RISC-V systems
        \item Investigating telemetry mechanisms based on hardware performance counters for heterogeneous RISC-V platforms
        \item Studying how performance monitoring data can inform scheduling and energy-aware optimisation
        \item Conducting experimental systems research on RISC-V platforms using Linux performance monitoring infrastructure
      \end{cvitems}
    }

  \cventry
    {Instituto Superior Técnico - Universidade de Lisboa}
    {MSc. Electrical \& Computer Engineering}
    {Lisbon, Portugal}
    {2021 -- 2023}
    {
      \begin{cvitems}
        \item{Thesis title: \cvwebsite{https://scholar.tecnico.ulisboa.pt/records/tn5twxt_kOmEMvcNvScCdHonti6KcnNRaybL}{Performance Monitoring and Event-based Sampling for RISC-V}}
        \item{Supervisors: Dr. Pedro \textsc{Tomás} and Dr. Nuno \textsc{Neves}}
        \item Developed performance monitoring tooling for RISC-V processors, focusing on hardware performance counters and profiling techniques
        \item Ported the PAPI (Performance API) library to the RISC-V architecture, enabling user-level access to hardware performance monitoring facilities
        \item Designed and implemented a proof-of-concept Precise Event Sampling (PES) mechanism for RISC-V processors on the CVA6 core
        \item Evaluated performance monitoring overheads and validated the implementation on a SiFive HiFive Unmatched platform
      \end{cvitems}
      \vspace{1em}
    }

  \cventry
    {Instituto Superior Técnico - Universidade de Lisboa}
    {BSc. Electrical \& Computer Engineering}
    {Lisbon, Portugal}
    {2017 -- 2021}
    {
      \begin{cvitems}
        \item Built a strong foundation in computer architecture, digital systems, and low-level programming, including processor organization, digital design, and algorithmic efficiency
        \item Developed analytical skills in advanced mathematics and computational methods, covering linear algebra, differential equations, probability, and numerical techniques for engineering systems
        \item Gained practical and theoretical experience in electronics, circuits, and instrumentation, including analog and digital electronics, signal measurement, and system characterization
        \item Studied signals, communications, and electromagnetic systems, including signal processing fundamentals and telecommunications principles
        \item Acquired multidisciplinary engineering knowledge in physics, control systems, and electrical energy systems, supporting a broad understanding of complex computing and electronic platforms
      \end{cvitems}
    }

\end{cventries}
